% SHARD-ID: LB-10-WAVE-OPERATORS
% SHARD-TITLE: Two-magnon wave operators, the Ward projection, and its erratum
% SHARD-SOURCES: theory/ml1-wave-operators.md; theory/ml4-ward-reduction.md;
%   theory/ml6-hs-sep.md; theory/ansatz-scattering-2m.md; definitions.md D31;
%   claims/CLAIMS.md rows ML1, ML1-D31-kernel, ML3, ML4, ML4-A, ML4-Ward,
%   ML6, ML6-HS-SEP, AC-EX-2M, AC-EX-2M-D29
% SHARD-CLAIMS: ML1, ML1-D31-kernel, ML4-Ward, ML4-A, ML4, ML3, ML6,
%   ML6-HS-SEP, AC-EX-2M, AC-EX-2M-D29, D31

\section{Two-magnon wave operators, the Ward projection, and its erratum}
\label{sec:wave-operators}

The previous section resolved the two-magnon Hilbert space spectrally. That is
not yet scattering theory: a Plancherel decomposition tells you what the
generalised eigenfunctions are, not that a wave packet prepared as two free
magnons in the far past actually converges to one of them. This section builds
the missing time-dependent layer, and then reports honestly on what happens
when one tries to turn a symmetry (Ward) identity into a soft theorem at finite
volume. The centre of gravity of the section is an erratum: a projected Ward
identity that looked exact turned out to hold at one overturned spin only. The
repair is displayed, the mechanism is named, and the downstream audit is
recorded. That kind of entry is what this labbook is for.

\subsection{The canonical two-magnon wave operators}

\begin{theorem}[Canonical two-magnon wave operators for the ferromagnet]
\label{thm:two-magnon-wave-operators}
\statusProved\ Assume \cref{def:fm-oracle-model,def:bethe-convention,def:bethe-soft-limit}
and $J>0$. In the centre/relative chart put
\begin{equation}
  \mathcal{C}:=(-\pi,\pi]\times(0,\pi),\qquad c(K):=\cos(K/2),\qquad
  k_1:=K/2-q,\quad k_2:=K/2+q,
\end{equation}
understood on the momentum torus, with $E(K,q)=J(2-2c(K)\cos q)$ and
$E_b(K)=J(1-c(K)^2)$. Let $\Sigma_{\mathrm{bad}}\subset\overline{\mathcal{C}}$
be the union of $q\in\{0,\pi\}$, of $K=\pi$, and of the curves on which either
$k_i\in\{0,\pi\}$; call a \emph{regular packet} any finite sum of smooth
compactly supported product packets whose support has positive distance from
$\Sigma_{\mathrm{bad}}$. Write
$\Hilb_{\mathrm{ch}}:=L^2(\mathcal{C},dK\,dq)$ and
$\Hilb_{b,\mathrm{ch}}:=L^2((-\pi,\pi],dK)$. Then there are canonical maps
$W_\pm:\Hilb_{\mathrm{ch}}\to\Hilb_2$ with the following properties.
\begin{enumerate}
\item $W_\pm$ are isometries with the same range $\Hilb_{\mathrm{sc}}$, and
  the two-magnon space splits orthogonally as
  $\Hilb_2=\Hilb_{\mathrm{sc}}\oplus\Hilb_b$ with
  $\Hilb_b\simeq\Hilb_{b,\mathrm{ch}}$.
\item They intertwine time and lattice translations, and
  \begin{equation}
    W_+^{*}W_-=M_{S_{\mathrm{phys}}},
    \label{eq:wpwm}
  \end{equation}
  multiplication by the physical outgoing/incoming ratio, which in this chart
  equals $S_{12}$ away from $K=\pi$ and has $\abs{S_{\mathrm{phys}}}=1$.
  Concretely $W_-=\mathcal{U}_{\mathrm{sc}}^{*}$ and
  $W_+=\mathcal{U}_{\mathrm{sc}}^{*}M_{S_{\mathrm{phys}}^{-1}}$, where
  $\mathcal{U}_{\mathrm{sc}}$ is the coefficient map against the
  delta-normalised scattering kernels \cref{eq:psi-sc}.
\item On every regular packet they are the strong time-dependent M\o{}ller
  limits
  \begin{equation}
    W_\pm F=\operatorname*{s-lim}_{t\to\pm\infty}e^{itH}\,I\,e^{-itE}F ,
    \label{eq:moller}
  \end{equation}
  for the canonical two-lowering identification $I$; so the limits exist and
  are isometric on precisely the packet class the campaign asked for.
\item In each relative-coordinate fibre with $K\ne0$ there is exactly one
  normalised bound eigenvector,
  $b_K(r)=\sqrt{1-c(K)^2}\,c(K)^{r-1}$ for $r\ge1$, of energy $E_b(K)$; at
  $K=0$ it degenerates into a threshold resonance and is not an $\ell^2$
  vector. No bound vector lies in either wave-operator range.
\item The one-magnon scattering space is
  $\Hilb_1=\ell^2(\Z)\simeq L^2(\mathbb{T},dk/2\pi)$; its incoming and
  outgoing maps coincide with the Fourier spectral map and are unitary.
\end{enumerate}
The union of regular packet spaces is dense in $\Hilb_{\mathrm{ch}}$, so the
compatible local isometries glue canonically to the global maps.
\end{theorem}
\begin{scope}
The time-dependent assertion is made only on regular packets. The extension
across $\Sigma_{\mathrm{bad}}$ --- equal velocities, one-particle endpoints,
the $K=\pi$ line --- is abstract, obtained by density from the isometries; no
endpoint or equal-velocity time-dependent statement is claimed, and no
incoming bound-pair channel is constructed. The bound band is a separate
stable composite channel, not a two-free-magnon channel.
\end{scope}
\provenance{ML1}{\texttt{theory/ml1-wave-operators.md} steps (1.1)--(1.7):
canonical spaces (1.2), stationary isometries and multiplier (1.3),
identification with the M\o{}ller limits and dense gluing (1.4), bound
isolation and one-magnon channel (1.5)}%
{\texttt{theory/checks/ml1\_ml6\_check.py} --- which checks only the explicit
fibre arithmetic, \emph{not} wave-operator existence}%
{\texttt{theory/verdicts/blitz-ml1-ml6-r1.md} (no fatal or major objection on
this row)}

\emph{Idea of proof.} The centre Fourier transform carries the chamber
$\{x<y\}$ to a direct integral of the half-line Jacobi fibres
\cref{eq:jacobi-fibre}, which differ from the free constant-coefficient
half-line operator only in the first diagonal entry. The Plancherel resolution
\cref{eq:plancherel} of \cref{thm:two-magnon-completeness} makes the pair of
coefficient maps against the scattering and bound kernels a unitary onto
$\Hilb_{\mathrm{ch}}\oplus\Hilb_{b,\mathrm{ch}}$ conjugating $H$ to
multiplication by $E$ and $E_b$; the two summands are therefore closed,
orthogonal and reducing, and exhaust $\Hilb_2$. In the chart the velocity
difference is $v(k_2)-v(k_1)=2J\cos(K/2)\sin q$, positive off the null line
$K=\pi$, so the incoming/outgoing labelling of the two relative branches is
fixed almost everywhere, giving the stationary maps and \cref{eq:wpwm}. To
identify these stationary maps with the time-dependent limits
\cref{eq:moller}, cover the compact support of a regular packet by finitely
many product rectangles, on each of which the two velocity supports have a
positive separation and all symbols are smooth; on each rectangle the
fixed-packet Cook construction of \cref{thm:ac-ex-2m} applies and its limit
agrees with the stationary map. Overlapping rectangles therefore agree, the
finite sum reproduces the packet, and density of the regular class extends
everything uniquely.

\subsection{Generalised one-magnon kernels beyond the product vacuum}

The wave operators above live on the exactly solvable ferromagnet. To carry
the same language to a general matrix-product vacuum one needs an abstract
hypothesis package, which the campaign calls $D31$. It is used twice in this
section, so it is stated here in full.

\begin{definition}[Exact fixed-packet two-magnon data over one vacuum]
\label{def:d31}
Fix one translation-invariant injective canonical matrix-product tensor $A$
and its state and GNS triple $(\omega_A,\Hilb_A,\Omega_A)$. All constructions
use this one vacuum representation; no covariant family of vacua is required.
Assume:
\begin{enumerate}
\item \emph{Dynamics and grading.} A finite-range translation-invariant
  Hamiltonian generates the dynamics, is implemented in $\Hilb_A$ by commuting
  positive-energy time and lattice translations, and is normalised by
  $H\Omega_A=0$, $\ker H=\C\Omega_A$. A conserved circle charge grades the
  selected magnon by one unit, and the construction lies in its charge-two
  sector.
\item \emph{An exact band.} There is a finite-multiplicity scalar magnon band
  $\omega\in C^2(\mathbb{T})$, isolated within the charge-one sector on the
  packet neighbourhoods of clause 4, and a Gram-normalised
  translation-covariant \emph{exact} map
  $\Gamma_M:L^2(\mathbb{T};\C^m)\to\Hilb_A$ with
  $H\Gamma_M=\Gamma_M(\omega\otimes1_m)$. The multiplicity is fixed;
  matrix-valued band crossings are outside the definition. Exactness is a
  hypothesis, not a Rayleigh--Ritz conclusion.
\item \emph{Filtered creators.} For each packet window, a spacetime-Schwartz
  filter of compact energy and momentum support, applied to a
  charge-one local observable, gives uniformly almost-local creators of
  positive energy transfer, translation-covariant as
  $a_{i,b}(n)=\tau_n(a_{i,b}(0))$, and acting on the vacuum as
  $a_{i,b}(n)\Omega_A=\Gamma_M(\chi_ie_n\otimes e_b)$ with
  $e_n(k)=e^{-ikn}$.
\item \emph{Packets.} Packet amplitudes are finite sums of $C_c^\infty$
  products. The dispersion is $C^\infty$ on neighbourhoods $U_i$ of the
  compact supports $K_i$; the velocity sets $V_i=\omega'(K_i)$ obey
  $\dist(V_1,V_2)\ge\varepsilon_v>0$ and $\dist(V_i,\{0\})\ge\varepsilon_0>0$.
  This last nonzero-velocity clause belongs to the campaign's packet register
  but enters no estimate of the theorem below; only $\varepsilon_v$ does.
  Filters satisfy $\chi_i=1$ on $K_i$ and $\supp\chi_i\subset U_i$; companion
  filters satisfy $\tilde\chi_i=1$ on $\supp\chi_i$ and
  $\supp\tilde\chi_i\subset U_i$; and $h_i=\omega\tilde\chi_i\in C^\infty(\mathbb{T})$.
\item \emph{A spectral inventory with margins.} Supply an inventory
  $\mathcal{R}_{\mathrm{inel}}$ of every alternative propagating charge-two
  band tuple $r=(r_1,\dots,r_n)$ expressible from the stated exact band data,
  and set $\theta_r(K):=\inf\{\sum_j\omega_{r_j}(p_j):\sum_jp_j=K,\ \sum_jq_{r_j}=2\}$.
  With $I_2=\{(\omega(k_1)+\omega(k_2),k_1+k_2):k_i\in K_i\}$, require for
  every $r$ that $\omega(k_1)+\omega(k_2)+\eta_{\mathrm{inel}}<\theta_r(k_1+k_2)$
  for one $\eta_{\mathrm{inel}}>0$ and all $k_i\in K_i$; absent channels count
  as $+\infty$. Every charge-two bound band in the supplied exact data is
  isolated, $d_B:=\inf_{k_i\in K_i,j}\abs{\omega(k_1)+\omega(k_2)-E_{B,j}(k_1+k_2)}>0$,
  and $P_B$ denotes the sum of its joint spectral projections; bound bands are
  present but are not incoming or outgoing two-magnon channels.
\end{enumerate}
No clustering inequality is assumed. Injectivity of $A$ already implies, for
every $\tilde\lambda$ between the transfer-spectrum gap and $1$ and for local
$C,D$ separated by $d_{\mathrm{sep}}$,
\begin{equation}
  \abs{\omega_A(CD)-\omega_A(C)\,\omega_A(D)}
  \le C_{A,\tilde\lambda}\norm{C}\norm{D}\,\tilde\lambda^{\,d_{\mathrm{sep}}},
  \label{eq:d31c2}
\end{equation}
and, after uniformly almost-local truncation, the four-operator version used
below. Both are consequences, not hypotheses. The package is fixed-packet
only: a soft scale and its limit are no part of it.
\end{definition}
\provenance{D31 (H-ACE2M)}{\texttt{definitions.md} D31, appended from
\texttt{theory/ansatz-scattering-2m.md} §7 with no clause edited}%
{\texttt{theory/checks/ansatz\_scattering\_2m\_check.py} corroborates only the
clustering identity \cref{eq:d31c2} on AKLT}%
{\texttt{theory/verdicts/ansatz-scattering-2m-r6.md} (the merge that gave the
clause set a numbered home)}

\begin{theorem}[Rigged generalised one-magnon kernels]
\label{thm:d31-kernel}
\statusSketch\ Assume \cref{def:d31} with multiplicity $m$, and rig the
one-magnon space as
$C^\infty(\mathbb{T};\C^m)\subset L^2(\mathbb{T},dk/2\pi;\C^m)\subset\mathcal{D}'(\mathbb{T};\C^m)$.
Transporting this rigging through the Gram-normalised exact band isometry
$\Gamma_M$ is intended to produce generalised kernels $\ket{k,b}$,
$1\le b\le m$, unique up to a measurable unitary change of multiplicity frame,
with
\begin{equation}
  \Gamma_Mf=\sum_b\int\frac{dk}{2\pi}f_b(k)\ket{k,b},
  \qquad
  \braket{k,b\,|\,k',b'}=2\pi\delta(k-k')\delta_{bb'},
  \qquad
  H\ket{k,b}=\omega(k)\ket{k,b},
\end{equation}
and with packets built from these kernels being exactly the one-particle legs
carried by the two-magnon wave operators, at the unit external-leg weight the
campaign's amplitude convention stipulates.
\end{theorem}
\begin{scope}
Two named gaps stand between this and a theorem, and both were raised by the
critic.
\begin{itemize}
\item \emph{Hypothesis binding.} The step identifying the kernels with the
  filtered creator legs uses clauses 3 and 4 of \cref{def:d31}, whereas the
  argument's own \textsc{assume} binds only clauses 1 and 2. Either the full
  package must be assumed, or the theorem must be split into a
  clause-1-and-2 part and a creator-leg part.
\item \emph{The rigging is the wrong one.} The global $C^\infty$ rigging is
  not preserved by multiplication by the dispersion, which \cref{def:d31} only
  requires to be $C^2$. The fix is to rig by $C^2$ instead, or to use a local
  rigging.
\end{itemize}
Either way the statement is conditional on \cref{def:d31}, and it asserts
nothing whatsoever about a merely variational excitation band.
\end{scope}
\provenance{ML1-D31-kernel}{\texttt{theory/ml1-wave-operators.md} step (1.6);
the defective leaves are the creator-leg identification and the choice of
rigging}{---}{\texttt{theory/verdicts/blitz-ml1-ml6-r1.md} (two major
objections, both named above)}

\subsection{The finite-sector Ward projection, and its erratum}

We now turn to the symmetry side. On a ring, in the sector $\Hilb_{n,N}$ with
$n$ overturned spins, write $S^\pm$ for the global spin raising and lowering
operators, $J^a_0$ for the three zero-momentum current components,
$D_{n,N}:=Q_0|_{\Hilb_{n,N}}$ for the restricted charge, and
$A_n:=D_{n,N}^{\dagger}D_{n,N}$. On a subspace where $A_n$ is invertible put
\begin{equation}
  P_{n,N}:=D_{n,N}A_n^{-1}D_{n,N}^{\dagger},
  \qquad
  R_{n,N}:=(1-P_{n,N})J^-_0 ,
  \label{eq:descendant-projector}
\end{equation}
so that $P_{n,N}$ is the orthogonal projection onto the whole descendant
subspace and $R_{n,N}$ is the part of the current orthogonal to it.

\begin{theorem}[Exact Ward projection in every finite sector]
\label{thm:ml4-ward}
\statusProved\ On $\Hilb_{n,N}$,
\begin{equation}
  D_{n,N}^{\dagger}J^-_0=2J^z_0+J^-_0S^+ .
  \label{eq:ward-8}
\end{equation}
On the highest-weight subspace $\ker S^+$ with $n<N/2$ the norm identity
\begin{equation}
  D_{n,N}^{\dagger}D_{n,N}=(N-2n)\,\mathbf 1
  \qquad\text{holds on }\ker S^+ ,
  \label{eq:ward-norm}
\end{equation}
and on any subspace where $A_n$ is invertible the general polar formula is
\begin{equation}
  P_{n,N}J^-_0=D_{n,N}A_n^{-1}\bigl(2J^z_0+J^-_0S^+\bigr).
  \label{eq:ward-polar}
\end{equation}
For one hard magnon of nonzero ring momentum $h$ this gives
\begin{equation}
  P_{1,N}J^-_0\ket{h}_N=\frac{2i\,v(h)}{N-2}\,Q_0\ket{h}_N ,
  \label{eq:ward-onehard}
\end{equation}
and the same identity holds diagonally for every hard packet supported away
from $h=0$.
\end{theorem}
\begin{scope}
These four statements are exact and were independently re-verified after the
erratum below. What they do \emph{not} give is a soft theorem: the
complementary vector $(1-P_{1,N})J^-_0\ket{h}_N$ is nonzero for generic $h$,
so the zero-momentum current is not velocity times total charge, and the Ward
identity computes the leading projection only.
\end{scope}
\provenance{ML4-Ward}{\texttt{theory/ml4-ward-reduction.md} step (1.3)
together with its erratum note}%
{\texttt{theory/checks/ml4\_check.py} (Ward and projection residuals at $n=1$:
maxima $8.899\cdot10^{-16}$ and $3.473\cdot10^{-15}$);
\texttt{theory/checks/ml4\_ward\_n2\_check.py} (the $n\ge2$ refutation and the
corrected form, red-capable)}%
{\texttt{theory/verdicts/corpus-r2.md} adjudication, scoped per
\texttt{theory/verdicts/soft-index-b-r1.md} F1 and the erratum work item;
audit \texttt{theory/verdicts/ml4-ward-n2-audit.md}}

\emph{Idea of proof.} The current is an $SU(2)$ vector, so
$[S^+,J^-_0]=2J^z_0$; expanding $S^+J^-_0$ with this commutator gives
\cref{eq:ward-8}. If $S^+\psi=0$ in the $n$-magnon sector then
$S^+S^-\psi=[S^+,S^-]\psi=2S^z\psi=(N-2n)\psi$, which is
\cref{eq:ward-norm}. Substituting \cref{eq:ward-8} into the definition
\cref{eq:descendant-projector} of the projector gives \cref{eq:ward-polar}.
For $n=1$ a single Fourier state is highest weight and satisfies
$J^z_0\ket{h}_N=iv(h)\ket{h}_N$ by direct evaluation, which yields
\cref{eq:ward-onehard}; translation covariance promotes it to packets.

\begin{refutedclaim}[The scalar projected Ward identity, refuted for two or more magnons]
\label{ref:ward-scalar}
\statusRefuted\ The withdrawn statement is the scalar display
\begin{equation}
  P_{n,N}J^-_0
  \overset{\text{\emph{false for }} n\ge2}{=}
  \frac{2}{N-2n}\,Q_0J^z_0
  \qquad\text{on }\ker S^+ .
  \label{eq:ward-refuted}
\end{equation}
It is exact at $n=1$ and false for every $n\ge2$.

\emph{Mechanism.} The derivation passes from the correct polar formula
\cref{eq:ward-polar} to \cref{eq:ward-refuted} by applying $A_n^{-1}$ as the
scalar $(N-2n)^{-1}$ to the vector $2J^z_0\psi$. That is legitimate only if
$2J^z_0\psi$ lies in $\ker S^+$, which is the subspace on which
\cref{eq:ward-norm} makes $A_n$ scalar. But $[S^+,J^z_0]=-J^+_0$, so
$S^+J^z_0\psi=-J^+_0\psi$, which vanishes at $n=1$ by momentum conservation
and is \emph{nonzero} for $n\ge2$: numerically $\norm{J^+_0\psi}=0$, $0.66$,
$1.36$ at $n=1,2,3$ for $N=8$. The scalar display then fails by
$0.26$ and $1.11$ at $n=2,3$ against left-hand sides of norm $0.91$ and
$1.39$, whereas the corrected form below is satisfied to $2.4\cdot10^{-15}$.

\emph{Surviving statement, and it is register-dependent.} Exactly at every
$n<N/2$, on $\ker S^+$,
\begin{equation}
  \boxed{\;P_{n,N}J^-_0=2\,D_{n,N}A_n^{-1}J^z_0\;}
  \label{eq:ward-corrected}
\end{equation}
--- that is, \cref{eq:ward-polar} with the $J^-_0S^+$ term dropped, because it
dies on $\ker S^+$, but with $A_n^{-1}$ kept as an \emph{operator} inverse,
$A_n$ being taken on the \emph{full} sector $\Hilb_{n,N}$, where it is not
scalar. The register matters and must always be named: in the
highest-weight-restricted register $D_\lambda:=Q_0|_{\ker S^+}$ one has
$A_\lambda=(N-2n)\mathbf 1$, so the same string $2D_\lambda A_\lambda^{-1}J^z_0$
collapses straight back to the refuted display. In that register the correct
form is instead
\begin{equation}
  P_{n,N}J^-_0=\frac{1}{m_\lambda}\,Q_0\,\Pi_{\mathrm{hw}}\,J^z_0,
  \qquad m_\lambda=\frac{N-2n}{2},
\end{equation}
with $\Pi_{\mathrm{hw}}$ the projection onto the highest-weight subspace.

\emph{A second casualty.} The two-hard specialisation, formerly written
$P_{2,N}J^-_0=\frac{2}{N-4}Q_0J^z_0$, is refuted by $100\%$ of its left-hand
side on the singular contact vector \cref{eq:singular-vector} of
\cref{thm:two-magnon-completeness} --- which is itself highest weight, since
applying $S^+$ to the alternating contact sum gives
$(-1)^x+(-1)^{x-1}=0$ at each site. At $N=8$ the true norm is $0.775$, the old
right-hand side has norm $1.414$, and the error is $0.775$ (relative error
$1.000$); at $N=10$ the error is $0.504$ (relative $0.730$). The corrected
full-sector form \cref{eq:ward-corrected} is exact there to
$2.4\cdot10^{-16}$. No scalar $(N-4)^{-1}$ version of the two-hard Ward part
is available.

\emph{Downstream audit.} The audit is complete: $44$ sites were examined, of
which $21$ were damaged, $2$ repairable by substitution, and $21$ clean, with
$0$ left unclear. \textbf{No headline claim was reached by the damage.}
Equations \cref{eq:ward-8,eq:ward-norm,eq:ward-polar,eq:ward-onehard} were
independently re-verified exact. The damaged statement sites in the two
earlier soft-index shards fold into a unified rewrite rather than being
patched twice.
\end{refutedclaim}
\provenance{ML4-Ward (erratum clause)}{\texttt{theory/ml4-ward-reduction.md},
the erratum notes at step (1.3) and at step (1.5).(2).1}%
{\texttt{theory/checks/ml4\_ward\_n2\_check.py} (green exit 0, \texttt{--red}
exit 1); \texttt{theory/checks/ml4\_check.py} at $n=1$}%
{\texttt{theory/verdicts/ml4-ward-n2-audit.md} (the complete 44-site audit);
first found in \texttt{theory/verdicts/soft-index-b-r1.md} F1}

\subsection{What actually kills the linear term}

The erratum has a constructive moral, isolated in the following abstract
lemma: the vanishing of the unwanted first-order contribution is not supplied
by the Ward identity at all. The Ward identity identifies the projected part
of the current. The zero comes from channel matching plus one derivative of
regularity.

\begin{lemma}[Matching plus $C^1$ regularity supplies the second soft zero]
\label{thm:ml4-a}
\statusProved\ Let $U:\Hilb\to\Kk$ be an isometry, $P=UU^{\dagger}$, and let
$\Gamma:(-\epsilon,\epsilon)\to B(\Hilb,\Kk)$ satisfy
\begin{equation}
  \Gamma(0)=U,\qquad
  \sup_{\abs{k}<\epsilon}\norm{\partial_k\Gamma(k)}\le C_\Gamma .
\end{equation}
Let $J\in B(\Hilb,\Kk)$, put $R=(1-P)J$, and let $b\in C^1$ satisfy $b(0)=0$
and $\abs{b(k)}\le C_b\abs{k}$. Then
\begin{equation}
  \norm{b(k)\Gamma(k)^{\dagger}R}\le C_bC_\Gamma\norm{J}\,k^2 ,
  \label{eq:ml4a-op}
\end{equation}
and for a normalised rescaled soft packet
$f_\varepsilon(k)=\varepsilon^{-1/2}f(k/\varepsilon)$,
\begin{equation}
  \norm{b(k)\Gamma(k)^{\dagger}Rg\,f_\varepsilon(k)}_{L^2(dk;\Hilb)}
  \le C_bC_\Gamma\norm{J}\,\varepsilon^2
    \norm{u^2f(u)}_{L^2(du)}\norm{g} .
  \label{eq:ml4a-packet}
\end{equation}
The same statement is allowed for channel traces rather than bounded maps,
provided the derivative bound is read in the trace norm on the chosen packet
class.
\end{lemma}
\provenance{ML4-A}{\texttt{theory/ml4-ward-reduction.md} step
(1.2)}{\texttt{theory/checks/ml4\_check.py} fixed-volume
probes}{\texttt{theory/verdicts/corpus-r2.md} adjudication}

\emph{Idea of proof.} Because $U$ is an isometry, $U^{\dagger}R=U^{\dagger}(1-UU^{\dagger})J=0$,
so $\Gamma(k)^{\dagger}R=[\Gamma(k)-\Gamma(0)]^{\dagger}R$. The fundamental
theorem of calculus bounds the difference by $C_\Gamma\abs{k}$, the projection
bounds $\norm{R}$ by $\norm{J}$, and multiplying by $\abs{b(k)}\le C_b\abs{k}$
gives \cref{eq:ml4a-op}; substituting $k=\varepsilon u$ and using that the
normalisation factor cancels the Jacobian gives \cref{eq:ml4a-packet}.

\begin{remark}[Where the zero comes from]
\label{rem:ml4a-moral}
The pair of properties that kills the orthogonal linear contribution is
\emph{energy-shell channel matching}, $\Gamma(0)=U$, together with \emph{$C^1$
on-shell trace regularity}. Energy conservation and multiplicity one imply the
first in a genuine two-body channel. The velocity sign $\sgn(v_h-v_s)$ only
labels which continuous wave is called outgoing; it supplies no zero at all.
\end{remark}

\subsection{Off-shell interpolation at fixed volume --- and why it is not a ring soft limit}

\begin{theorem}[Fixed-volume off-shell interpolation for one hard magnon]
\label{thm:ml4-offshell}
\statusSketch\ Fix $I=[a,b]\Subset(0,\pi)$ and let $\Gamma_N(k)$ be the
analytic interpolation obtained by transporting the coordinate scattering
formula from total momentum $h+k$ to the $k=0$ fibre, normalised by
$\Gamma_N(k)\ket{h}:=B_N(k,h)/\sqrt{N-2}$, where on $0\le x<y<N$
\begin{equation}
  B_N(k,h;x,y)=\frac{e^{ih(x+y)/2}}{\sqrt N}
    \bigl\{s(k,h)e^{iq(y-x)}+e^{-iq(y-x)}\bigr\},
  \qquad q=\tfrac{h-k}{2},
\end{equation}
with $s(k,h)$ the solution of the contact equation \cref{eq:contact-equation}
analytic at zero. Let
\begin{equation}
  \mathsf{A}_{\perp,N}(k;f,g):=(e^{ik}-1)\,
   \braket{\Gamma_N(k)f,\;R_{n,N}\,g}
\end{equation}
be the interpolated orthogonal contribution to the Ward-reduced connected
numerator. At zero soft momentum the interpolation matches the descendant,
$\Gamma_N(0)\ket{h}=Q_0\ket{h}_N/\sqrt{N-2}$, since $s(0,h)=1$. For every
fixed $N$ there is then a constant $C'_{I,N}<\infty$ with
\begin{equation}
  \abs{\mathsf{A}_{\perp,N}(k;f,g)}
   \le C'_{I,N}\,k^2\,\norm{f}_{I,N}\norm{g}_{I,N}
\end{equation}
for arbitrary smooth compactly supported hard packets; the corresponding
soft-packet norm is accordingly $O_N(\varepsilon^2)$, and consequently
\begin{equation}
  \lim_{N\to\infty}\ \lim_{k\to0}\
   \frac{\mathsf{A}_{\perp,N}(k;f,g)}{k}=0 ,
\end{equation}
where the inner limit is the \emph{off-shell analytic} limit.
\end{theorem}
\begin{scope}
This is emphatically not an on-shell ring soft limit, and the distinction is
not pedantic. At fixed periodic $N$ the nonzero momenta are discrete, so
$k\to0$ at fixed $N$ never runs through physical ring states; the interpolated
$\Gamma_N(k)$ is an off-shell analytic vector. Worse, the first \emph{physical}
sequence refutes volume uniformity outright. Take $h=2\pi/5$, ring sizes
divisible by five, and $k_N=2\pi/N$. The low-frequency relative-coordinate sum
obeys
\begin{equation}
  \lim_{N\to\infty}N^{-1}\braket{B_N(c/N,h),B_N(0,h)}
   =\frac{8\bigl[1-\cos(c/2)\bigr]}{c^2},
\end{equation}
which at $c=2\pi$ equals $4/\pi^2$ rather than the value $1$ at $c=0$. Hence
the transported trace tends to $2\abs{v(h)}(1-4/\pi^2)>0$ along that sequence,
so its ratio to $\abs{k_N}$ grows linearly in $N$ and the normalised version
grows like $\sqrt N$: no $N$-independent version of the estimate holds. What
remains open, therefore, is exactly the statement the campaign wants --- a
packet-smeared infinite-volume on-shell trace bound controlling the regime
$k=\Theta(1/N)$ --- together with the two-hard, three-magnon channel discussed
next.
\end{scope}
\provenance{ML4}{\texttt{theory/ml4-ward-reduction.md} step (1.4) together
with step (1.5).(2).2; narrowed after the erratum, since the earlier
step (1.5).(2).1 consumed the refuted scalar display}%
{\texttt{theory/checks/ml4\_check.py}: fixed-volume fitted exponents
$0.99700$--$0.99988$ for the raw orthogonal trace and
$1.99700$--$1.99988$ after multiplication by $e^{ik}-1$, with the omitted
projection as built-in red mutation; the joint-scaling probe at $k=2\pi/N$ and
its \texttt{--red-uniform} mutation reinstating the retracted uniform claim,
which exits nonzero}{\texttt{theory/verdicts/corpus-r2.md} (held at sketch)}

\begin{remark}[The sharp obstruction in the two-hard channel]
\label{rem:ml4-3}
To infer a quadratic bound for the orthogonal part in the two-hard sector one
would need a three-magnon outgoing channel satisfying the two hypotheses of
\cref{rem:ml4a-moral} in a packet trace norm uniform in volume.
\cref{thm:two-magnon-completeness} resolves the \emph{input} two-magnon state
and says nothing about the three-magnon target. Unlike the two-body half-line
fibre, a three-body fixed-energy shell can have several open channels; energy
conservation alone then does not force every zero-soft channel into the range
of the charge, and an orthogonal degenerate channel can contribute at first
order after the single lattice difference. Excluding that multiplicity is the
sharp missing hypothesis, and it needs three-body wave operators together with
a limiting-absorption or trace estimate. Assuming closed Bethe factorisation
would hide the issue rather than settle it, and is not permitted here.
\end{remark}

\begin{conjecture}[Uniform current form-factor regularity]
\label{conj:ml3}
\statusConjecture\ The finite-volume regularity of the current form factors
upgrades, after wave-packet smearing, to a bound uniform in volume and in the
soft limit. Any proof must control the regime $k=\Theta(1/N)$ and must exclude
a physical $1/k_s$ current pole, showing that the delta and principal-value
pieces appearing in the coefficient formulas are distributional normalisation
artefacts.
\end{conjecture}
\provenance{ML3}{---}{---}{--- (future work)}

\begin{conjecture}[Order of limits and channel contamination in general]
\label{conj:ml6}
\statusConjecture\ The order of the infinite-volume limit, the wave-packet
width, and the soft limit can be controlled in general, and the bound state
and off-shell coefficients do not contaminate the leading real scattering
channel.
\end{conjecture}
\provenance{ML6}{---}{---}{--- (future work)}

\subsection{The order of limits on a separated packet class}

The general conjecture above is open. On a restricted but honest class of
preparation protocols it is a theorem, and the following is what the campaign
actually owns about limit order.

\begin{theorem}[Ordered limits remove bound and string weight on the separated class]
\label{thm:ml6-hs-sep}
\statusProved\ Assume \cref{def:fm-oracle-model,def:bethe-convention,def:bethe-soft-limit},
$J>0$, and a fixed hard interval $I=[a,b]\Subset(0,\pi)$. Take a one-sided
normalised soft packet $f_\varepsilon(k)=\varepsilon^{-1/2}f(k/\varepsilon)$
with $\supp f\subset[c_1,c_2]\subset(0,\infty)$ and a normalised hard packet
$g_\sigma$ supported in $I$. For each fixed $\varepsilon>0$ use the separated
preparation sequence: translate the incoming hard packet by $R_j\to\infty$,
take a settling time $T_j$ with $d_\varepsilon T_j-R_j\to\infty$, impose the
separation condition, and take the ring size large enough that the
thermodynamic and sampling error tends to zero. The order is
\begin{equation}
  N\to\infty\ \prec\ (R,T,\sigma)_j\to(\infty,\infty,0)\ \prec\ \varepsilon\downarrow0 .
  \label{eq:l6-order}
\end{equation}
Then:
\begin{enumerate}
\item at each fixed $\varepsilon$, the total squared coefficient of the
  prepared state outside the selected incoming real-scattering packet tends to
  zero --- this includes the two-string and bound band and all
  real-scattering coefficients off the selected support --- and the analogous
  outgoing statement holds after the settling time;
\item along every actual row-measure limit subsequence the normalised
  connected readout is exactly
  \begin{equation}
    \mathsf{A}_*(\varepsilon)
     =\int\bigl[S_{\mathrm{phys}}(k,h)-1\bigr]\,d\mu_{*,\varepsilon}(k,h);
    \label{eq:hs-sep-readout}
  \end{equation}
\item with $\bar k_*(\varepsilon):=\int k\,d\mu_{*,\varepsilon}$,
  \begin{equation}
    \mathsf{A}_*(\varepsilon)=2i\,\bar k_*(\varepsilon)+O_I(\varepsilon^2),
  \end{equation}
  the remainder being uniform in the hard packet and in the actual limit
  measure, so neither bound nor off-shell coefficients contaminate the leading
  real scattering channel;
\item no soft-uniform Cook estimate is needed, because the fixed-$\varepsilon$
  channel limit is completed first.
\end{enumerate}
\end{theorem}
\begin{scope}
The theorem proves no uniform pointwise limit of individual coefficients ---
only that their packet-smeared contaminating sum is controlled. It applies to
the separated primitive packet class, not to an arbitrary preparation
protocol. It constructs no incoming bound-pair wave operator, makes no
endpoint or equal-velocity claim, and has no extension to $a\downarrow0$: that
is the same non-uniformity as $k_h\to0$ recorded in
\cref{rem:not-a-scattering-length}, and it is excluded rather than overcome.
\end{scope}
\provenance{ML6-HS-SEP}{\texttt{theory/ml6-hs-sep.md} steps (1.1)--(1.9):
spectral fence (1.2), collective coefficient estimate (1.3), ordered limit
(1.4), soft expansion in the purified channel (1.5), endpoint fence (1.6)}%
{\texttt{theory/checks/ml1\_ml6\_check.py}: green exit 0, with
\texttt{--red-bound} and \texttt{--red-unitarity} exiting 1; arithmetic
only}{\texttt{theory/verdicts/blitz-ml1-ml6-r1.md}}

\emph{Idea of proof.} At the same total momentum the continuum-to-bound
separation is $E(K,q)-E_b(K)=J\abs{e^{iq}-c}^2$, which at $k=0$ equals
$J\sin^2(h/2)$, so compactness of $I$ gives a uniform positive fence. The
overlap of the unseparated charge-created relative wave with the bound fibre
is exactly
\begin{equation}
  \beta(K,q)=\frac{2\sqrt{1-c^2}\,(\cos q-c)}{1-2c\cos q+c^2},
  \qquad
  \beta\bigl(h+k,\tfrac{h-k}{2}\bigr)=2k+O_I(k^2),
\end{equation}
so the raw bound amplitude of a soft packet is $O_I(\varepsilon)$ --- the same
order as the connected law itself, hence \emph{not} by itself enough to
exclude contamination. The decisive input is the ordered limit. At fixed
$\varepsilon$, the separated preparation makes the incoming Cook tail decay
like a high inverse power of the separation, so the projection onto the bound
band and onto real roots off the selected window vanishes in norm
\emph{before} the soft limit is taken; Parseval on the ring, together with
strong convergence of the finite fibre matrices to the half-line matrices,
converts this into a collective statement about smeared coefficients, with no
root-by-root labelling and no principal-value cancellation required. Only then
is $\varepsilon\downarrow0$ taken, at which point the expansion
\cref{eq:S-phys} of \cref{thm:two-body-soft} is integrated against the limit
measure. Because the contaminating projection is already exactly zero, it
cannot contribute a competing term of order $\varepsilon$.

\subsection{Cook limits over a general matrix-product vacuum}

\begin{theorem}[Fixed-packet two-magnon Cook limits from exact band data]
\label{thm:ac-ex-2m}
\statusProvedConditional\ Assume \cref{def:d31}. Put
$\Hilb_{0,12}:=L^2(K_1;\C^m)\otimes L^2(K_2;\C^m)$ with free Hamiltonian
$H_{0,12}=\omega\otimes1+1\otimes\omega$ on the smooth product core
$\mathcal{D}_{12}$, the two slots labelling the two disjoint packet windows
rather than distinguishable species, and let $I$ be the two-creator
identification map. Then, \emph{as a conditional implication from
\cref{def:d31} and nothing else}:
\begin{enumerate}
\item[(i)] the fixed-packet limits
  $W_\pm F:=\lim_{t\to\pm\infty}e^{itH}Ie^{-itH_{0,12}}F$ exist for
  $F\in\mathcal{D}_{12}$;
\item[(ii)] $W_\pm$ are isometries, extend uniquely to the packet-domain
  closure, and intertwine joint time and space translations; their ranges are
  annihilated by the sum $P_B$ of the listed fibrewise-isolated charge-two
  bound-band projections, and are orthogonal to every alternative propagating
  channel in the inventory, whose joint energies exceed
  $\sup_{I_2}E+\eta_{\mathrm{inel}}$.
\end{enumerate}
Specialised to the ferromagnet of \cref{def:fm-oracle-model}, a fixed packet
range lies in the matching part of the scattering summand of
\cref{thm:two-magnon-completeness} and is orthogonal to its two-string
summand, and $W_+^{*}W_-$ is exactly multiplication by the physical ratio
$S_{\mathrm{phys}}$ of \cref{def:bethe-convention} (the inverse ratio under
the opposite velocity ordering), of modulus one.

The clustering estimates \cref{eq:d31c2} and its four-operator version are
\emph{derived} from injectivity of the vacuum tensor, not assumed; no
independent clustering hypothesis enters.

The implication is not vacuous: the ferromagnet of
\cref{def:fm-oracle-model} satisfies all five clauses of \cref{def:d31}. Its
all-up state is the injective bond-dimension-one product tensor, its shifted
Hamiltonian is positive with $\ker H=\C\Omega$, the total magnetisation
supplies the conserved circle grading, the Fourier map
$g\mapsto\int\frac{dk}{2\pi}g(k)\sum_xe^{ikx}\ket{x}$ is an exact isometric
band map, a spacetime-Schwartz filter of $S^-_n$ gives the covariant creators,
and \cref{thm:two-magnon-completeness} supplies exactly one bound band with a
positive gap and an empty inelastic inventory.
\end{theorem}
\begin{scope}
Clauses 1--5 of \cref{def:d31} --- the exact band, the creators, the
covariance and the threshold data --- are \emph{hypotheses}, not consequences
of a variational ansatz; only the clustering estimates are derived. Nothing
here is uniform in a soft scale: every estimate is at fixed packets, and the
Cook and Gram constants carry the packet Schwartz seminorms $s_N(F)$. Along a
soft rescaling $f_\varepsilon(k)=\varepsilon^{-1/2}f(k/\varepsilon)$ one has
$\norm{\partial^jf_\varepsilon}_\infty=\varepsilon^{-j-1/2}\norm{f^{(j)}}_\infty$,
so the $N$ integrations by parts carry a factor $\varepsilon^{-N}$ and the
bound degrades without limit as $\varepsilon\downarrow0$. The nonzero-velocity
constant $\varepsilon_0$ of clause 4 enters no estimate in this theorem; only
the velocity separation $\varepsilon_v$ does. The theorem asserts nothing
about the fixed-time interface of \cref{thm:d29-interface}, which no step of
its proof consumes. No endpoint or equal-velocity construction, no
soft-uniform Cook bound, no bound-state wave operator, no compatibility or
range exhaustion across packet windows, and no asymptotic completeness is
claimed.
\end{scope}
\provenance{AC-EX-2M (parts A2M.1--A2M.2)}%
{\texttt{theory/ansatz-scattering-2m.md} steps (1.2)--(1.7$'$): fixed-packet
fence (1.2), clustering upgrade (1.3), filtered one-particle equations (1.4),
Cook integrability (1.5), isometry (1.6), intertwining and the exact match
(1.7), nonvacuity on the ferromagnet (1.7$'$)}%
{\texttt{theory/checks/ansatz\_scattering\_2m\_check.py}: green exit 0 and
seven registered red modes exiting 1 under optimised Python, all nine gates
evaluated before exit. The certificate corroborates \emph{only} the clustering
identity \cref{eq:d31c2} on AKLT, reproduced independently in a second tensor
basis; it tests no Cook limit, no exact band, no spectral separation, and no
identification with the ferromagnet's scattering summand. Two gate caveats are
on record: the rescaled agreement gate does not certify the absence of a
two-sided support-length factor (support-length independence is proved
analytically, not numerically), and one guard gate is identically zero by
algebra, so it is a code-shape guard rather than a numerical
certificate.}{\texttt{theory/verdicts/ansatz-scattering-2m-r5.md} §9}

\emph{Idea of proof.} Exactness of the band map means the filtered creator
defect $D_{i,b}(n)=[H,a_{i,b}(n)]-\sum_mh_i(m-n)a_{i,b}(m)$ annihilates the
vacuum and is almost local uniformly in $n$. Expanding $H$ on the two-creator
vector and cancelling both filtered free convolutions therefore leaves only a
commutator term,
$(HI-IH_{0,12})F=\sum_{x,y}F(x,y)[D_1(x),a_2(y)]\Omega_A$, whose norm decays
faster than any power of $x-y$. Outside its velocity cone each free packet is
rapidly decreasing in $\abs{x}+\abs{t}$, while on the product of the two main
cones the velocity separation forces $\abs{x-y}\ge\varepsilon_v\abs{t}/2$;
splitting the sum accordingly gives an integrable majorant of order
$\abs{t}^{-3}$, and the fundamental theorem of calculus then gives both strong
limits. Isometry follows by applying the derived clustering estimate to the
four-creator scalar product on the product cones, where the Gram-normalised
exact band maps factorise into the free tensor Gram form, with the complement
again power-suppressed. Intertwining is the usual replacement of $t$ by $t+s$;
the joint spectral calculus then converts the fibrewise bound gap and the
inelastic margin into exact annihilation of the corresponding projections.
Finally, on the ferromagnet, stationary phase puts the two relative branches
of \cref{eq:psi-sc} in the incoming and outgoing regions exactly as the
coefficient convention of \cref{def:bethe-convention} prescribes, so the
constructed labels are that convention's and not a new one.

\subsection{The fixed-time interface: what is and is not established}

\begin{theorem}[The fixed-time protocol interface]
\label{thm:d29-interface}
\statusSketch\ For the adjudicated fixed-time interface, two \emph{separate}
statements are established, and they are not composed.
\begin{enumerate}
\item \emph{Diagonal compactness, conditional on assumed bounds.} Write the
  amputated datum at an allowed full index tuple $\alpha$ as
  $\mathsf{A}_\alpha(h)=\mathsf{N}_\alpha(h)/\mathsf{D}_\alpha(h)$ for $h\in I$.
  If the selected-hard-packet bounds
  \begin{equation}
    \operatorname*{ess\,inf}_\alpha\ \operatorname*{ess\,inf}_{h\in I}
      \abs{\mathsf{D}_\alpha(h)}\ge d_I>0,
    \qquad
    \sup_\alpha\norm{\mathsf{N}_\alpha}_{L^2(I)}\le C_I<\infty
    \label{eq:d29-den}
  \end{equation}
  hold for every allowed tuple, then $\sup_\alpha\norm{\mathsf{A}_\alpha}_{L^2(I)}\le C_I/d_I$,
  so every full-index sequence respecting the precedence
  \begin{equation}
    N\to\infty\ \prec\ t\to\pm\infty\ \prec\ (W\uparrow\Z,\sigma\downarrow0)_j
     \ \prec\ \varepsilon\downarrow0
    \label{eq:d29-order}
  \end{equation}
  with $\varepsilon\downarrow0$ has a weakly convergent subsequence in
  $L^2(I)$. This is Banach--Alaoglu applied to assumed bounds: both displayed
  inequalities in \cref{eq:d29-den} are assumptions, not derivations, and no
  iterated-order datum is asserted to exist.
\item \emph{Creator-choice independence --- a theorem about Haag--Ruelle
  families.} At fixed $\varepsilon>0$, a soft creator family that is
  admissibly Haag--Ruelle, satisfies the one-particle on-shell normalisation,
  and has velocity support disjoint from the hard packet, may be replaced by
  the filtered creator of \cref{def:d31} without changing any connected
  on-shell pairing in the limit of large times.
\end{enumerate}
\end{theorem}
\begin{scope}
\textbf{The second statement does not apply to the adjudicated interface.}
Both protocols on record specify a \emph{fixed-time} insertion
$Q[f_\varepsilon]\psi$ of the charge on an already prepared hard vector, and
that is not a Haag--Ruelle creator family: it never free-evolves its profile.
The failure is precisely of Haag--Ruelle admissibility and not of the
one-particle normalisation, which \emph{is} proved on the ferromagnet for the
family built from that same charge acting on the vacuum.

For the adjudicated datum the exact mismatch is visible at the soft law's own
linear order. By \cref{thm:two-body-soft} and \cref{eq:R8-mismatch},
\begin{equation}
  Q_{k_s}\ket{k_h}-\ket{B^{\mathrm{in}}}
   =(1-S_{12})\ket{P_{12}}
   =-2ik_s\ket{P_{12}}+O(k_s^2),
\end{equation}
a relative size $\sqrt2\,k_s(1+O(k_s))$ --- nonzero exactly where the soft law
lives. The identification of the fixed-time datum with the constructed channel
of \cref{thm:ac-ex-2m} is therefore \textbf{open}, and two steps are named as
missing:
\begin{itemize}
\item[(i)] prove, at order $k_s$ and with a displayed uniform remainder, that
  the large-time readout's connected packet-amputated on-shell pairing equals
  the constructed-channel one \emph{despite} that mismatch;
\item[(ii)] exhibit an instance of the soft-regularity condition for the
  fixed-time family, including existence of its volume and time limits.
\end{itemize}
The soft-regularity clauses have been verified for exactly one datum: the
\emph{constructed-channel} family on the infinite chain, which has a single
member and involves no exhaustion,
\begin{equation}
  \mathsf{A}(\varepsilon)(h)
   =\int d\mu_f(u)\,S_{\mathrm{phys}}(\varepsilon u,h)
\end{equation}
on the ferromagnet, for which every uniformity clause is trivially satisfied
because the family is a singleton and the limit-existence clause is bypassed
rather than met. It therefore backs no claim about the interface. The
window-and-width-uniform form of the soft-regularity condition is open on
every model, the ferromagnet included; it remains a hypothesis with no
verified instance. No fixed-time charge-to-channel identification, no
fixed-time charge-to-scattering-vector equality, and no first-jet theorem for
the interface is claimed.
\end{scope}
\provenance{AC-EX-2M-D29}{\texttt{theory/ansatz-scattering-2m.md} steps
(1.8)--(1.9), excluding step (1.9).(2).4, which belongs to
\cref{thm:ac-ex-2m}; the ported creator-choice theorem is taken from
\texttt{refs/arxiv-1412.2970} (final clause of its Haag--Ruelle
theorem)}{No certificate exists for this row:
\texttt{theory/checks/ansatz\_scattering\_2m\_check.py} tests none of the
assumed bounds \cref{eq:d29-den}, Haag--Ruelle admissibility, the missing step
(i), the volume and time limits, or the soft-regularity
condition}{\texttt{theory/verdicts/ansatz-scattering-2m-r5.md} §9 and
\texttt{ansatz-scattering-2m-r6.md}}
